% Part: first-order-logic
% Chapter: introduction
% Section: satisfaction

\documentclass[../../../include/open-logic-section]{subfiles}

\begin{document}

\olfileid{fol}{int}{sat}

\olsection{સંતોષ}

!!{formula}s શું છે તે જાણ્યા પછી આપણે તરત પ્રથમ-ક્રમ તર્કશાસ્ત્રના
અર્થવિચાર તરફ આગળ વધી શકીએ: અહીં પાયાની વ્યાખ્યા !!a{structure}ની
છે. આપણી સાદી ભાષા માટે !!a{structure}~$\Struct M$ના માત્ર ત્રણ ઘટકો
છે: \emph{!!{domain}} કહેવાતો અરિક્ત ગણ~$\Domain{M}$,
~$\Struct M$માં $\Obj a$ જેને નિર્દેશે છે તે વસ્તુ અને~$\Struct M$માં
$\Obj P$ જે વસ્તુઓ માટે સત્ય છે તે વસ્તુઓ. $\Obj a$થી નિર્દેશાતી
વસ્તુને~$\Assign{\Obj a}{M}$ અને $\Obj P$ જેના માટે સત્ય છે એ
વસ્તુઓના ગણને~$\Assign{\Obj P}{M}$થી દર્શાવીએ છીએ.
!!^a{structure}~$\Struct{M}$ બરાબર આ ત્રણ વસ્તુઓની બનેલી છે:
$\Domain{M}$, $\Assign{\Obj a}{M} \in \Domain{M}$ અને
$\Assign{\Obj P}{M} \subseteq \Domain{M}$. સામાન્ય કિસ્સો વધુ જટિલ
હશે, કારણ કે તેમાં ઘણાં !!{predicate}s અને !!{constant}s હશે,
!!{predicate}s એક કરતાં વધુ સ્થાનોવાળાં હોઈ શકે અને !!{function}s પણ
હશે.
\sourcecorrection{OLFINT-004}{મૂળની \texttt{satisfaction.tex},
પંક્તિ~22માં એક કરતાં વધુ સ્થાનોવાળાં ``અચળો'' કહેવામાં આવ્યા છે.
અચળો એક વસ્તુને નિર્દેશે છે; સામાન્ય ભાષામાં અનેક સ્થાનોવાળાં
વિધેયો આવે છે. તેથી અહીં ``વિધેયો'' પુનઃસ્થાપિત કર્યું છે.}

!!{variable}s વિનાનાં !!{formula}s માટે સંતોષ વ્યાખ્યાયિત કરવા આટલું
પૂરતું છે. આપણે !!{formula}sને જે રીતે વ્યાખ્યાયિત કર્યાં છે તેનું
પ્રતિબિંબ પાડતી અનુમાનાત્મક વ્યાખ્યા આપવાનો વિચાર છે. આણ્વિક સૂત્ર
~$\Struct{M}$માં ક્યારે સંતોષાય તે સ્પષ્ટ કરીએ અને પછી, દાખલા તરીકે,
$!A$,~$\Struct{M}$માં સંતોષાય છે કે નહિ તેના આધારે $\lnot !A$,
~$\Struct{M}$માં ક્યારે સંતોષાય તે કહીએ. દાખલા તરીકે, આ રીતે
વ્યાખ્યાયિત કરી શકીએ:
\begin{enumerate}
  \item $\Atom{\Obj P}{\Obj a}$,~$\Struct{M}$માં ત્યારે અને માત્ર
  ત્યારે સંતોષાય જ્યારે $\Assign{\Obj a}{M} \in \Assign{\Obj P}{M}$.
  \item $\lnot !A$,~$\Struct{M}$માં ત્યારે અને માત્ર ત્યારે સંતોષાય
  જ્યારે $!A$,~$\Struct{M}$માં સંતોષાતું ન હોય.
  \item $(!A \land !B)$,~$\Struct{M}$માં ત્યારે અને માત્ર ત્યારે
  સંતોષાય જ્યારે $!A$,~$\Struct{M}$માં સંતોષાય અને $!B$ પણ
  ~$\Struct{M}$માં સંતોષાય.
\end{enumerate}
ધારો કે $\Domain{M} = \{0, 1, 2\}$, $\Assign{\Obj a}{M} = 1$ અને
$\Assign{\Obj P}{M} = \{1, 2\}$. આ વ્યાખ્યા આપણને કહેશે કે
$\Atom{\Obj P}{\Obj a}$,~$\Struct{M}$માં સંતોષાય છે (કારણ કે
$\Assign{\Obj a}{M} = 1 \in \{1,2\} = \Assign{\Obj P}{M}$). તે એ પણ
કહે છે કે $\lnot \Atom{\Obj P}{\Obj a}$,~$\Struct{M}$માં સંતોષાતું
નથી, પરિણામે $\lnot\lnot \Atom{\Obj P}{\Obj a}$ સંતોષાય છે અને
$(\lnot \Atom{\Obj P}{\Obj a} \land \Atom{\Obj P}{\Obj a})$
સંતોષાતું નથી, અને આમ આગળ.

પરિમાણકો માટે વ્યાખ્યા આપવા જઈએ ત્યારે મુશ્કેલી આવે છે: આપણને
``$\lexists[\Obj v_0][\Atom{\Obj P}{\Obj v_0}]$ ત્યારે અને માત્ર
ત્યારે સંતોષાય જ્યારે $\Atom{\Obj P}{\Obj v_0}$ સંતોષાય'' જેવું કંઈક
કહેવું ગમશે. પરંતુ !!{structure}~$\Struct{M}$, !!{variable}sનું શું
કરવું તે કહેતું નથી. હકીકતમાં આપણે એવું કહેવા માગીએ છીએ કે
$\Atom{\Obj P}{\Obj v_0}$,~$\Obj v_0$ના \emph{કોઈ એક મૂલ્ય માટે}
સંતોષાય છે. આ વાત ચોક્કસ કરવા $\Obj a$ને જ નહિ પણ~$\Obj v_0$ને પણ
~$\Domain{M}$નાં !!{element}s નિયુક્ત કરવાની રીત જોઈએ. આ માટે આપણે
!!{variable}-\emph{નિયુક્તિઓ} દાખલ કરીએ છીએ. !!^a{variable}-નિયુક્તિ
એટલે માત્ર એવું વિધેય~$s$ જે !!{variable}sને~$\Domain{M}$નાં
!!{element}s સાથે જોડે (આપણા દાખલામાં $0$, $1$ અથવા~$2$માંથી કોઈ એક
સાથે). કોઈ !!a{formula}માં કયાં !!{variable}s આવશે એ અગાઉથી જાણતા નથી,
તેથી $s$ કયાં !!{variable}sને મૂલ્ય આપે તેના પર મર્યાદા મૂકી શકતા નથી.
સરળ ઉકેલ એ છે કે $s$ બધાં !!{variable}s~$\Obj v_0$, $\Obj v_1$,
\dots{}ને મૂલ્યો આપે એવી માગણી કરીએ\@. આપણને જરૂરી હોય તે જ મૂલ્યો
વાપરીશું.
\sourcecorrection{OLFINT-005}{મૂળની \texttt{satisfaction.tex},
પંક્તિ~65માં ચલનાં સંભવિત મૂલ્યો $1$, $2$, $3$ લખ્યાં છે, પરંતુ
નિશ્ચિત વ્યાપકક્ષેત્ર $\{0,1,2\}$ છે. અહીં તેની સાથે સુસંગત મૂલ્યો
$0$, $1$, $2$ રાખ્યાં છે.}

!!{formula}sનો સંતોષ માત્ર !!a{structure} સાપેક્ષ વ્યાખ્યાયિત કરવાને
બદલે આપણે તેને !!a{structure}~$\Struct{M}$ \emph{અને} !!a{variable}
નિયુક્તિ~$s$ સાપેક્ષ વ્યાખ્યાયિત કરીશું અને ટૂંકમાં
$\Sat{M}{!A}[s]$ લખીશું. !!{variable}s ધરાવતાં આણ્વિક !!{formula}s
સંભાળવા હવે આપણી વ્યાખ્યામાં વધારાનો કિસ્સો આવશે:
\begin{enumerate}
  \item $\Sat{M}{\Atom{\Obj P}{\Obj a}}[s]$ ત્યારે અને માત્ર ત્યારે
  જ્યારે $\Assign{\Obj a}{M} \in \Assign{\Obj P}{M}$.
  \item $\Sat{M}{\Atom{\Obj P}{\Obj v_i}}[s]$ ત્યારે અને માત્ર ત્યારે
  જ્યારે $s(\Obj v_i) \in \Assign{\Obj P}{M}$.
  \item $\Sat{M}{\lnot !A}[s]$ ત્યારે અને માત્ર ત્યારે જ્યારે
  $\Sat{M}{!A}[s]$ નહિ.
  \item $\Sat{M}{(!A \land !B)}[s]$ ત્યારે અને માત્ર ત્યારે જ્યારે
  $\Sat{M}{!A}[s]$ અને $\Sat{M}{!B}[s]$ બંને.
\end{enumerate}
હવે એક મુશ્કેલી ઉકેલાઈ: $s(\Obj v_0)$ મૂલ્ય માટે~$\Struct{M}$,
$\Atom{\Obj P}{\Obj v_0}$ને ક્યારે સંતોષે છે તે કહી શકીએ છીએ.
$\lexists[\Obj v_0][\Atom{\Obj P}{\Obj v_0}]$ માટે વ્યાખ્યા બરાબર
બેસાડવા એક વાત વધુ કરવી પડશે: આપણે એવું જોઈએ છીએ કે
$\Sat{M}{\lexists[\Obj v_0][\Atom{\Obj P}{\Obj v_0}]}[s]$ ત્યારે અને
માત્ર ત્યારે જ્યારે $\Obj v_0$ને મૂલ્ય નિયુક્ત કરવાની \emph{કોઈ એક}
રીત~$s'$ માટે $\Sat{M}{\Atom{\Obj P}{\Obj v_0}}[s']$ હોય. પરંતુ
$\Obj v_0$ને નિયુક્ત કરેલું મૂલ્ય $s(\Obj v_0)$ જે વસ્તુને નિર્દેશે છે
તે જ હોય એવું જરૂરી નથી. તેના માટે સંજ્ઞા દાખલ કરીએ: જો $m \in
\Domain{M}$ હોય, તો $\Subst{s}{m}{\Obj v_0}$ એવી નિયુક્તિ છે જે
~$s$ જેવી જ છે ($\Obj v_0$ સિવાયનાં બધાં !!{variable}s માટે), પરંતુ
$\Obj v_0$ને~$m$ નિયુક્ત કરે છે. હવે આપણી વ્યાખ્યા આ થઈ શકે:
\begin{enumerate}\setcounter{enumi}{4}
  \item $\Sat{M}{\lexists[\Obj v_i][!A]}[s]$ ત્યારે અને માત્ર ત્યારે
  જ્યારે કોઈ~$m \in \Domain{M}$ માટે
  $\Sat{M}{!A}[\Subst{s}{m}{\Obj v_i}]$.
\end{enumerate}
શું આ કામ કરે છે? ધારો કે દરેક~$i \in \Nat$ માટે $s(\Obj v_i) = 0$
રાખીએ. $\Sat{M}{\lexists[\Obj v_0][\Atom{\Obj P}{\Obj v_0}]}[s]$
ત્યારે અને માત્ર ત્યારે જ્યારે એવો કોઈ $m \in \Domain{M}$ હોય કે
$\Sat{M}{\Atom{\Obj P}{\Obj v_0}}[\Subst{s}{m}{\Obj v_0}]$. આવું
મૂલ્ય છે: આપણે $m = 1$ અથવા $m = 2$ લઈ શકીએ. નોંધો કે $s$ પોતે
$\Obj v_0$ને નિયુક્ત કરે છે એ મૂલ્ય~$s(\Obj v_0)$---અહીં $0$---કામ ન
કરે તોપણ આ સત્ય છે. આપણને $\Sat{M}{\Atom{\Obj P}{\Obj
v_0}}[\Subst{s}{1}{\Obj v_0}]$ મળે છે, પરંતુ
$\Sat{M}{\Atom{\Obj P}{\Obj v_0}}[s]$ મળતું નથી.

આ ગૂંચવણભર્યું અને ભારે લાગતું હોય, તો એ ખરેખર એવું છે. પરંતુ બધાં
!!{formula}s માટે સંતોષની ચોક્કસ, અનુમાનાત્મક વ્યાખ્યા આપવા વધારાની આ
જટિલતા જરૂરી છે અને અર્થવિચારી ખ્યાલોને ચોક્કસ રીતે વ્યાખ્યાયિત કરવા
આવું કંઈક જોઈએ જ. આ કરવાની બીજી રીતો પણ છે, પણ બધી સરખી જ
(અ)સુઘડ છે.

\end{document}
